\documentclass[11pt]{article}
\usepackage[margin=1in]{geometry}
\usepackage{fontspec}
\usepackage{titlesec}
\usepackage{enumitem}
\usepackage{longtable}
\usepackage{array}
\usepackage{booktabs}
\usepackage{xcolor}
\usepackage{hyperref}
\usepackage{fancyhdr}
\usepackage{tcolorbox}
\usepackage{amssymb}
\usepackage{amsmath}
\usepackage[normalem]{ulem}
\usepackage{listings}
\usepackage{graphicx}

\lstset{
  basicstyle=\fontsize{6pt}{7pt}\ttfamily,
  breaklines=true,
  breakatwhitespace=false,
  columns=fullflexible,
  frame=single,
  framerule=0.3pt,
  rulecolor=\color{gray!40},
}
\tcbuselibrary{skins,breakable}

\definecolor{marlinblue}{HTML}{1B4965}
\definecolor{marlinteal}{HTML}{2E8B8B}
\definecolor{lightgray}{HTML}{F4F4F4}
\definecolor{codebg}{HTML}{F7F7F9}
\definecolor{risk}{HTML}{B23A48}

\titleformat{\section}{\large\bfseries\color{marlinblue}}{\thesection}{0.5em}{}[{\color{marlinblue}\titlerule}]
\titleformat{\subsection}{\normalsize\bfseries\color{marlinteal}}{\thesubsection}{0.5em}{}

\hypersetup{colorlinks=true, linkcolor=marlinblue, urlcolor=marlinteal}

\pagestyle{fancy}
\fancyhf{}
\fancyhead[L]{\small\textcolor{gray}{Marlin FD + PETSc Roadmap}}
\fancyhead[R]{\small\textcolor{gray}{\thepage}}
\renewcommand{\headrulewidth}{0.3pt}

\setlength{\parskip}{0.5em}
\setlength{\parindent}{0em}
\sloppy
\emergencystretch=3em
\Urlmuskip=0mu plus 1mu

\newenvironment{schematic}{%
  \begin{tcolorbox}[breakable, colback=codebg, colframe=gray!40, boxrule=0.4pt, arc=1mm, left=4mm, right=4mm, top=2mm, bottom=2mm]
  \ttfamily\footnotesize
}{%
  \end{tcolorbox}
}

\newcommand{\pr}[2]{\subsection{PR#1 --- #2}}
\newcommand{\status}[1]{\textcolor{marlinteal}{\textbf{Status:}} #1}

\title{\textcolor{marlinblue}{\Huge Marlin Finite-Difference + PETSc\\[2pt] Implicit Solve}\\[8pt]
\Large Development Roadmap}
\author{}
\date{\today}

\begin{document}
\maketitle

\begin{tcolorbox}[colback=lightgray, colframe=marlinblue, boxrule=0.6pt, arc=1mm]
\textbf{Status of this document: living reference.} Update it as each PR
lands or as decisions change. Each PR section carries a status
(\texttt{proposed} / \texttt{approved} / \texttt{implemented} / \texttt{verified}) ---
keep that current rather than treating this document as a one-time plan.
\end{tcolorbox}

\tableofcontents
\vspace{1em}

\section{Context and Goal}

Marlin composes PDE physics as \texttt{TensorCompute} objects
(\texttt{include/tensor\_computes/TensorOperatorBase.h}), today either
FFT-spectral (\texttt{ReciprocalMatDiffusion}, \texttt{ReciprocalAllenCahn}, \ldots)
or one real-space FD operator (\texttt{FiniteDifferenceLaplacian}) stepped
purely explicitly (\texttt{RealSpaceForwardEuler}). There is no PETSc usage in
Marlin's own tensor solve path, and no true boundary-condition mechanism for
the real-space path --- smooth boundary method (SBM) no-flux today is faked
via a multiplicative \texttt{psi} mask on the RHS, not a real boundary
stencil.

\textbf{Goal:} compose PDE terms as TensorComputes (plus a new BC sub-block of
TensorComputes), recognize which terms are linear-with-known-stencil, and
solve the implied $Ax=b$ for those via PETSc Krylov solvers applied
\textit{matrix-free} (operator action computed on the fly from the same
stencil coefficients already used for explicit evaluation --- see
Section~\ref{sec:partial-assembly} --- not via an assembled sparse matrix),
with options configured MOOSE-style in the solve block; lag everything else
into $b$; optionally Picard-iterate; allow non-time-dependent fields (e.g.\ a
Poisson equation) to be solved alongside transient ones
(Section~\ref{sec:steady-mixing}). Target ``basic'' functionality: spinodal
decomposition (Cahn-Hilliard), KKS/grand-potential, and the smooth boundary
method --- all currently periodic-only and spectral.

This plan was produced by: (a) exploring the existing TensorCompute /
TensorSolver / TensorProblem / DomainAction architecture, MOOSE's
\texttt{PetscSupport} options-DB mechanism, and MOOSE's \texttt{FixedPointSolve}
Picard knobs; (b) surveying existing example input files for the target
physics (\texttt{test/tests/kks/KKS\_no\_flux\_bc.i},
\texttt{test/tests/cahnhilliard/*.i}, \texttt{examples/cahn\_hilliard/*.i}
including the FEM comparison variant \texttt{cahnhilliard\_fem.i},
\texttt{examples/libtorch\_kks/KKS\_libtorch.i}); (c) a design pass; (d)
decisions with the stakeholder on four consequential questions
(Section~\ref{sec:decisions}); (e) a critical review surfacing residual
risks (Section~\ref{sec:risks}) that this roadmap is structured to retire
early and cheaply, rather than discovering them deep into a large PR.

\section{Target Architecture}
\label{sec:architecture}

This section is a reference for what we are building toward across the PR
sequence in Section~\ref{sec:roadmap} --- no single PR implements all of it
at once.

\subsection{Hybrid Design: PETSc Owns the Algorithms, Torch Owns the Kernels}
\label{sec:hybrid}

\begin{tcolorbox}[colback=lightgray, colframe=marlinblue, boxrule=0.6pt, arc=1mm]
\textbf{Decided 2026-10-18} after the PR1a/PR1b evidence chain
(Section~\ref{sec:pr1b}). The deciding factor was generality and long-term
maintenance (PETSc's solver stack is maintained by full-time scientists),
not speed: the measured gap was not large in either direction.
\end{tcolorbox}

\begin{longtable}{p{0.36\textwidth} p{0.1\textwidth} p{0.44\textwidth}}
\toprule
Layer & Owner & Mechanism \\
\midrule
\endhead
Krylov method, convergence tests & PETSc & \texttt{KSPCG} (SPD), \texttt{KSPGMRES}/\texttt{KSPFGMRES} otherwise \\
Multigrid cycle and coarse solve & PETSc & \texttt{PCMG}; coarse level is a small assembled \texttt{Mat} \\
Default smoother & PETSc & \texttt{KSPCHEBYSHEV} + \texttt{PCJACOBI} per level \\
Later: SNES, TS, AMG & PETSc & AMG on an assembled \texttt{Pmat} (PR4a) \\
Level operator action and diagonal & torch & \texttt{MATSHELL} \texttt{MatMult}/\texttt{GetDiagonal} running the \texttt{LinearFDOperator} stencil \\
Restriction / prolongation & torch & \texttt{MATSHELL}s via \texttt{PCMGSetRestriction}/\texttt{PCMGSetInterpolation} \\
Optional custom smoother & torch & \texttt{PCSHELL} per level (e.g.\ symmetric red-black Gauss-Seidel for non-SPD), optionally fused via \texttt{TracedComputeSequence} \\
\bottomrule
\end{longtable}

Configuration rules from the evidence:
\begin{itemize}[leftmargin=*]
\item \textbf{Mark SPD operators.} Set \texttt{MAT\_SPD} on every level
whose operator is SPD. This halved PCMG setup on the H200 (724$\to$408~ms
at $n=8192$, 69$\to$9~ms at $n=256$) with the same iterations and solution.
\item \textbf{Leave a hook for analytic Chebyshev bounds.} Gershgorin bounds
from the stencil coefficients would be passed through
\texttt{KSPChebyshevSetEigenvalues} and would remove the remaining per-level
eigenvalue estimate. Not built in v1.
\item \textbf{Reuse the preconditioner across Picard iterations}
(Section~\ref{sec:decisions}, item~3).
\end{itemize}

The torch-native solvers remain. \texttt{conjugateGradientSolve} is the MPS
path, and the PR1b torch multigrid spike is the reference for a fused,
torch-only cycle.

\subsection{Operator Application: Matrix-Free (Partial Assembly)}
\label{sec:partial-assembly}

\begin{tcolorbox}[colback=lightgray, colframe=marlinblue, boxrule=0.6pt, arc=1mm]
\textbf{This design uses matrix-free operator application by default} ---
the MFEM-style ``partial assembly'' model, as opposed to MOOSE's
full-Jacobian-assembly model where PDE terms are assembled into a sparse
\texttt{Mat}.
\end{tcolorbox}

Concretely: $A$'s action on a vector is computed on the fly from the same
stencil coefficients \texttt{FiniteDifferenceLaplacian} and
\texttt{FiniteDifferenceDiffusion} already use for their matrix-free
\texttt{computeBuffer()} (a \texttt{torch::nn::functional::convNd} call);
no sparse PETSc \texttt{Mat} is written out. PETSc's KSP is handed
this operator via \texttt{MatCreateShell} (its \texttt{MatMult} callback
invokes the existing conv-kernel code), so Krylov methods (CG, GMRES) work
unmodified, with no sparse matrix ever stored or assembled.
Preconditioning uses matrix-free options (Jacobi/Chebyshev, or geometric
multigrid built from the same stencil action at coarser grid levels) rather
than algebraic multigrid, since AMG (\texttt{-pc\_type gamg}/\texttt{hypre})
needs the assembled matrix it is designed to precondition.

This reverses the originally-proposed assembled-\texttt{Mat} default, now
kept only as a documented fallback for cases this matrix-free approach
can't handle well (see Section~\ref{sec:soa}, PR4a) because: (a) Marlin already
has the matrix-free kernels and a structured grid, which is exactly the
regime partial assembly is built for; (b) it avoids a round-trip through
PETSc's own \texttt{Vec}/\texttt{Mat} memory layout on every solve, which
matters doubly on Marlin's GPU-resident (CUDA/MPS) torch tensors; (c) PETSc
has no Metal/MPS GPU backend at all (CUDA/HIP/Kokkos only), so on an
MPS-only machine a torch-native matrix-free action (which \textit{can} run
on MPS) is a structural win over anything routed through PETSc's own device
vectors; (d) the matrix-free action operates directly on each field's whole
dense grid tensor (structure-of-arrays, Section~\ref{sec:soa}), the layout
that vectorizes under SIMD on CPU and coalesces on GPU --- the same property
that motivates (a)--(c) independent of which device is in use. See the new
PR1a benchmarking task (Section~\ref{sec:roadmap}) for actually measuring
torch-vs-PETSc overhead before committing further.

\subsection{Data Layout: Structure-of-Arrays, and the Zero-Copy Requirement at the PETSc Boundary}
\label{sec:soa}

Every \texttt{TensorCompute} buffer is one dense grid-shaped
\texttt{torch::Tensor} per field: structure-of-arrays (SoA). This is a
different layout than MOOSE's per-DOF array-of-structures (AoS), where
coupled variables are interleaved per node for sparse assembly.
\texttt{computeBuffer()} and the
matrix-free stencil action (Section~\ref{sec:partial-assembly}) operate
elementwise/stencil-wise over a whole contiguous per-field array, which is
the layout SIMD vectorization and GPU coalesced access both need. The
segregated-per-field-solve decision (Section~\ref{sec:decisions}, item 1)
and the per-field \texttt{TimeDerivativeOperator} contribution
(Section~\ref{sec:steady-mixing}) both preserve this: each PETSc object the
plan builds is single-DOF-per-grid-point for its field, with no
interleaving. PR8 (deferred, Section~\ref{sec:roadmap}), if ever built,
must use field-split/block \texttt{Vec}s rather than interleaved per-point
DOFs, so a future monolithic solver does not reintroduce AoS layout by a
different route.

\begin{tcolorbox}[colback=lightgray, colframe=marlinblue, boxrule=0.6pt, arc=1mm]
\textbf{Zero-copy is a requirement at the point where a field's torch
tensor becomes a PETSc \texttt{Vec}.}
\end{tcolorbox}

PETSc supports wrapping an existing buffer instead of allocating and copying
into its own:
\begin{itemize}
\item \textbf{CPU}: \texttt{VecCreateSeqWithArray}/\texttt{VecCreateMPIWithArray}
  wrap the host buffer directly --- no allocation, no copy, provided the
  tensor is contiguous and the local chunk matches PETSc's partition.
\item \textbf{CUDA}: \texttt{VecCreateSeqCUDAWithArrays}/
  \texttt{VecCreateMPICUDAWithArrays} (and \texttt{VecCUDAPlaceArray} to
  swap a device pointer into an existing \texttt{Vec}) wrap a raw CUDA
  device pointer --- \texttt{tensor.data\_ptr<double>()} on a CUDA tensor is
  such a pointer, so a torch-CUDA field can be handed to PETSc with no host
  round-trip and no device-to-device copy, subject to the same
  partition-matching requirement as R1, now at the pointer level rather
  than logical indexing.
\item \textbf{MPS}: no equivalent exists --- PETSc has no Metal/MPS
  \texttt{Vec} type at all, so there is no pointer to wrap. The only options
  on an MPS tensor are a host copy (pay the transfer every solve) or skip
  PETSc on that device entirely. The torch-native fallback (PR1a) is
  accordingly the only zero-copy-capable path on MPS.
\end{itemize}

Zero-copy memory sharing and zero-overhead compute are different claims.
Wrapping avoids the copy, but PETSc's own \texttt{Vec} arithmetic inside
\texttt{KSPSolve} (dot products, norms, AXPY) still runs through PETSc's
device backend (cuBLAS/cuSPARSE for \texttt{VECCUDA}), a different kernel
implementation than libTorch's own CUDA kernels; the matrix-free
\texttt{MatMult} callback is the one point execution crosses back into
torch's kernel per iteration. This plan claims zero-copy memory sharing; it
makes no claim about matching compute kernels inside \texttt{KSPSolve}
itself.

The zero-copy requirement covers the operator matrix (\texttt{Amat},
the \texttt{MatCreateShell} the Krylov method multiplies against every
iteration) and the solution/RHS \texttt{Vec}s. \texttt{KSPSetOperators(ksp,
Amat, Pmat)} accepts a separate \texttt{Pmat} used only to build the
preconditioner, which can be
a real assembled sparse \texttt{Mat} fed to \texttt{PCHYPRE}/BoomerAMG or
\texttt{PCGAMG} without requiring \texttt{Amat} itself to be assembled
(PR4a). \texttt{Pmat} need not be an exact Jacobian --- AMG preconditioners
work from an approximate or frozen-coefficient matrix as long as it
captures the right coarsening/anisotropy structure. This is a scoped,
infrequent exception (one sparse assembly at the reassembly cadence of
Section~\ref{sec:decisions}, item 3, not a per-iteration cost) to the
zero-copy guarantee, not a hole in it.

\subsection{Class Schematic}

\begin{schematic}
\begin{verbatim}
                MooseObject, DependencyResolverInterface, MarlinConstantInterface
                                      |
                             TensorOperatorBase
                        (computeBuffer, realSpaceComputeBuffer,
                         getInputBuffer/getOutputBuffer, updateDependencies)
                                      |
                +---------------------+----------------------+
                |                                            |
         TensorOperator<T>                              TensorSolver
      (adds "buffer" output param)                 (substep() loop driver)
                |                                            |
 +--------------+---------------+          +------------------+------------------+
 |                               |          |                                     |
FiniteDifferenceLaplacian  FiniteDifferenceDiffusion  PetscFDSolverBase    PicardTensorSolver
 (existing; + mixin)         (new; variable-coeff)   (new: ONE solver       (new: wraps a
    |                               |                  object; owns DM,    PetscFDSolverBase
    +---------------+---------------+                  loops one-by-one    as _compute,
                     |                                  over tracked       Picard-iterates it)
                     |                                  fields, each with
                     |                                  a MatShell (matrix-
                     |                                  free action) + KSP,
                     |                                  no assembled sparse
                     |                                  Mat by default
             LinearFDOperator   <--- mixin interface, NOT a TensorOperatorBase
         (applyAt / matrix-free action,              subclass -- multiply-inherited
          rhsContributionAt,                          alongside TensorOperator<T>
          coefficientsChanged,
          optional stencilAt() ONLY for the
          documented assembled-Mat fallback)

       TensorOperatorBase
               |
      TensorBoundaryCondition  (new; "variable", "boundary", a/b/c params)
               |
         TensorRobinBC  (new; general a*u + b*du/dn = c, both fillGhostCells
               |         and stencilContribution derived from the same a/b/c)
    +----------+----------+
    |                     |
TensorDirichletBC    TensorNeumannBC   (thin presets: a=1,b=0  /  a=0,b=1)
\end{verbatim}
\end{schematic}

\subsection{Data-Flow Schematic --- One Picard Iteration}

\texttt{PetscFDSolverBase} is a single solver object (not one object per
field). Its \texttt{substep()} simply loops over its list of tracked fields,
one at a time, in dependency order, doing one assemble+solve per field
before moving to the next:

\begin{schematic}
\begin{verbatim}
                 +-----------------------------------------------+
                 |  for each tracked field, one-by-one, in        |
                 |  dependency order via DependencyResolverInterface |
                 +-----------------------------------------------+
                                    |
  +---------------------------------+----------------------------------+
  v                                 v                                  v
bind MatShell(field)        assembleRHS(field)                 (BCs consulted by both)
  MatShell's MatMult callback  walk owned DMDA points            TensorBoundaryCondition
  invokes this field's          for each LinearFDOperator         ::applyGhostContribution()
  LinearFDOperator(s) matrix-     ::rhsContributionAt()             -- modifies the matrix-free
  free action directly on       for each lagged nonlinear op:       action/RHS near physical
  torch tensors (same convNd      computeBuffer() (matrix-free,     (non-periodic) domain edges
  kernel as the existing           evaluated at current/lagged
  matrix-free compute path)        buffer state) -> added to b
  -- no MatSetValues, no sparse
  Mat ever assembled
  |                                 |
  +----------------+----------------+
                    v
              KSPSolve(ksp, b, x)
       (Krylov method + matrix-free
        preconditioner -- Jacobi/Chebyshev
        or geometric MG using the same
        stencil action; KSP configured via
        its own -<field>_ petsc_options
        prefix, set MOOSE-style in the
        solve block)
                    |
                    v
        copy Vec x -> torch buffer for field
                    |
                    v
 TensorProblem::exchangeGhostLayers(field)
(periodic MPI halo exchange  +  TensorBoundaryCondition
 ::fillGhostCells() for physical, non-periodic edges --
 SAME a/b/c as stencilContribution() above, by construction)
                    |
                    v
    (back to outer loop: next field, same iteration)

After all fields updated once:
    PicardTensorSolver checks convergence
    (residual norm or named Postprocessor,
     picard_rel_tol/abs_tol, picard_max_its)
    -- not converged -> apply relaxation_factor,
       re-invoke PetscFDSolverBase::substep()
    -- converged / max_its reached -> done, advance timestep
\end{verbatim}
\end{schematic}

\subsection{Boundary-Condition Dual-Use Invariant}

This is the design detail most likely to silently rot if not actively
enforced by a test --- flagged explicitly as Risk R4 in
Section~\ref{sec:risks}.

\begin{schematic}
\begin{verbatim}
TensorRobinBC(a, b, c) for buffer "u" on face "left"
     |
     +--> fillGhostCells(u)            [explicit / matrix-free path]
     |      ghost value = one-sided extrapolation from a,b,c
     |      consumed by: nonlinear computeBuffer() evaluations,
     |                   explicit solvers, visualization
     |
     +--> applyGhostContribution(idx)  [implicit / matrix-free path]
            action/RHS edit derived from the SAME a,b,c
            consumed by: PetscFDSolverBase's MatMult callback / assembleRHS
            (an assembled-matrix stencilAt() variant of this exists only
             for the PR4a fallback)

INVARIANT: both derive from one a,b,c triple in one class.
REQUIRED TEST (not optional): a test that runs BOTH paths on the same
BC instance and checks they imply the same boundary value/flux --
not just that each individually "looks reasonable" in isolation.
\end{verbatim}
\end{schematic}

\subsection{Steady-State (Non-Time-Dependent) Fields Mixed with Transient Ones}
\label{sec:steady-mixing}

New requirement: it must be possible to solve a non-time-dependent PDE
(e.g.\ a Poisson equation for an electric/chemical potential)
\textit{alongside} transient phase-field variables in the same simulation,
with the steady field feeding source terms into the transient ones'
nonlinear remainder (and vice versa).

Design implication: the time-derivative term must stop being an implicit
assumption baked into \texttt{PetscFDSolverBase} and instead become
\textbf{just another optional \texttt{LinearFDOperator} contribution},
registered per field like any other. A transient field registers a
time-derivative operator (contributing $u/dt$ to the matrix-free action and
$u_{\text{old}}/dt$ to $b$); a steady field registers none --- its
matrix-free action is the bare steady operator (e.g.\ the Laplacian for
Poisson), and \texttt{PetscFDSolverBase::substep()} simply solves it to
convergence every time its lagged inputs change (every Picard iteration, or
every timestep if Picard is disabled), exactly like any other tracked field
in the one-by-one loop --- no special-casing needed in the solver's control
flow, only in what each field registers.

\begin{schematic}
\begin{verbatim}
PetscFDSolverBase, one-by-one field loop, THIS run has 3 tracked fields:

  field "c"   (transient, Cahn-Hilliard concentration)
     linear ops: [FiniteDifferenceDiffusion(D), TimeDerivative(dt)]
  field "eta" (transient, Allen-Cahn order parameter)
     linear ops: [FiniteDifferenceLaplacian(kappa), TimeDerivative(dt)]
  field "phi" (STEADY, e.g. electric potential)
     linear ops: [FiniteDifferenceLaplacian(epsilon)]   <- no TimeDerivative
     nonlinear remainder: charge density sourced from "c"/"eta" (lagged)

Picard loop still just re-invokes substep() over all 3 fields in
dependency order; "phi" re-solves to a fresh steady state every
iteration using the latest lagged c/eta, "c"/"eta" see the latest phi
as a lagged source in their own nonlinear remainder.
\end{verbatim}
\end{schematic}

This is additive to the architecture already described, not a redesign ---
flagged here because it changes \texttt{LinearFDOperator}'s contract (time
derivative is explicit, not assumed) and needs its own new class,
\textbf{\texttt{TimeDerivativeOperator}}
(\texttt{include/tensor\_computes/TimeDerivativeOperator.h/.C}, implementing
\texttt{LinearFDOperator}), plus its own verification case (PR5a).

\section{Decisions Already Made}
\label{sec:decisions}

\begin{enumerate}[leftmargin=*]
\item \textbf{Segregated per-field linear systems, solved one-by-one by a
single solver object} (\texttt{PetscFDSolverBase}), Picard-coupled across
fields --- not monolithic block coupling. One solver object looping over
fields (rather than one solver object per field) keeps the input-file
syntax and the Picard wrapper simple: Picard just re-invokes the same
\texttt{substep()}, which always walks every tracked field in sequence.
\textit{Future-proofing:} the \texttt{LinearFDOperator::StencilEntry} carries
an optional \texttt{target\_field} (unused in v1, defaults to ``self''), and
DOF-indexing is kept as its own component rather than hardwired
one-field-per-Mat, so a future monolithic solver is additive, not a rewrite.
\item \textbf{Mutually exclusive solve mode} --- a simulation picks PETSc-FD
mode or an existing explicit/spectral mode, not a mix in one run.
\item \textbf{Preconditioner reuse across Picard iterations, with an
iteration threshold and a mandatory rebuild-and-retry} (revised 2026-10-18;
replaces ``reassembly every Picard iteration'').
\begin{itemize}
\item \textbf{What is reused.} The shell operators always carry the current
coefficients, so the Krylov method solves the current system. Only the PCMG
setup is reused: the Jacobi diagonals, the Chebyshev estimates and the
coarse factorisation.
\item \textbf{When it is rebuilt.} When the iteration count exceeds a
threshold, and always, with a retry, on a negative converged reason. A stale
Chebyshev setup went indefinite (\texttt{KSP\_DIVERGED\_INDEFINITE\_PC}) in
PR1b.
\item \textbf{How a rebuild is triggered.} Bumping a shell Mat's state makes
\texttt{PCApply} lazily rerun \texttt{PCSetUp} inside the next solve.
\end{itemize}
Measured at $\beta=3$, $n=4096$: 1695~ms with reuse vs.\ 2786~ms rebuilding
every step. Reuse saves little at $\sim$1300$\times$ coefficient contrast.
\item \textbf{Dirichlet, Neumann, and Robin all in v1} --- Robin is the
general implementation; Dirichlet/Neumann are presets of it.
\item \textbf{Hybrid PETSc + torch} (revised 2026-10-18,
Section~\ref{sec:hybrid}), with KSPCG + PCMG on torch shell levels as the
default for SPD fields, and \textbf{matrix-free (partial-assembly) operator
application by default.} PETSc's KSP operates against a \texttt{MatCreateShell} whose
\texttt{MatMult} calls the existing matrix-free \texttt{computeBuffer()}-style
stencil action directly on torch tensors; no sparse matrix is assembled or
stored for \texttt{Amat}. For cases where matrix-free preconditioning
proves insufficient on anisotropic/heterogeneous coefficients, PR4a adds a
separately assembled \texttt{Pmat} to feed AMG
(\texttt{PCHYPRE}/\texttt{PCGAMG}) while \texttt{Amat} stays matrix-free
(Section~\ref{sec:soa}, R7); a fully assembled \texttt{Amat} is kept only as
a documented last-resort fallback beyond that. Rationale for the
matrix-free default: Marlin already has the matrix-free kernels and
structured grid partial assembly is built for; it keeps the
structure-of-arrays per-field layout (Section~\ref{sec:soa}), which
vectorizes under SIMD on CPU and coalesces on GPU; and PETSc has no
Metal/MPS GPU backend at all (CUDA/HIP/Kokkos only), so on an MPS-only
machine the torch-native matrix-free action has a structural advantage.
The \texttt{Vec}s PETSc's KSP operates on must wrap each field's existing
tensor buffer zero-copy rather than allocate and copy into PETSc's own
memory --- a requirement (Section~\ref{sec:soa}). On CPU and CUDA a
wrapping path exists
(\texttt{VecCreateMPIWithArray}/\texttt{VecCreateMPICUDAWithArrays}); on MPS
none exists, which is why MPS runs the torch-native fallback instead.
\item \textbf{Steady-state (non-time-dependent) fields may be freely mixed
with transient fields} in one simulation (e.g.\ a Poisson equation for a
potential, solved every Picard iteration/timestep alongside transient
Cahn-Hilliard/Allen-Cahn fields) --- see Section~\ref{sec:steady-mixing}.
This does \textit{not} reopen decision \#2 (explicit/spectral vs.\ PETSc-FD
solve modes stay mutually exclusive); it only means multiple fields
\textit{within} one PETSc-FD-mode simulation can differ in whether they
carry a time-derivative term.
\end{enumerate}

\section{Residual Risks This Roadmap Retires Early}
\label{sec:risks}

\begin{longtable}{>{\bfseries}p{0.06\textwidth} p{0.32\textwidth} p{0.3\textwidth} p{0.24\textwidth}}
\toprule
\# & Risk & Why it matters & Where it's retired \\
\midrule
\endhead
R1 & DMDA partition layout may not actually match Marlin's torch-tensor
layout rank-for-rank & Silent cross-rank data corruption if wrong, hard to
detect later & PR2, in isolation, before any solver logic is built on it \\
R2 & Picard may not converge (or converge slowly) for strongly nonlinear
coupling (e.g.\ KKS mobility) & Could force a redesign of the coupling
strategy late & PR6, explicitly measured and tuned on a real coupled case,
not assumed \\
R3 & PCMG setup may dominate Picard runtime, and reuse may degrade the
preconditioner at high contrast & Could make the ``performant'' goal fail
in practice & Partly measured in PR1b: setup is 1--1.8$\times$ a solve
after \texttt{MAT\_SPD}; reuse saved 39\% at $\beta=3$. Re-profiled at PR3
and PR6 \\
R4 & Ghost-fill and stencil-row BC math could silently diverge over time as
code evolves & Would reintroduce exactly the kind of ad hoc inconsistency
SBM's \texttt{psi} mask already has & A same-BC-instance dual-path test is
a \textit{required} deliverable of PR1, not optional polish \\
R5 & Robin generality could block Dirichlet/Neumann usability if sequenced
badly & Delays the SBM no-flux use case, which only needs Neumann & PR1
delivers Dirichlet+Neumann usable on their own; Robin's extra generality
lands in the same PR but is not a gate for anything downstream \\
R6 & Overhead of crossing between torch tensors and PETSc's own
\texttt{Vec}/\texttt{Mat}/device memory (especially on CUDA/MPS) could
erase any benefit of using PETSc at all, vs.\ a custom torch-native
matrix-free CG/multigrid solver, if zero-copy \texttt{Vec} wrapping
(Section~\ref{sec:soa}) is not actually achieved in practice & Could
invalidate the entire ``use PETSc'' premise of this plan on GPU & PR1a,
measured directly for both wrapped-zero-copy and naive-copy cases (CPU
here; CUDA run handed off to the user) before PR2/PR3 commit further
engineering to the PETSc path; PR3 re-confirms zero-copy directly, not by
timing inference \\
R7 & Matrix-free preconditioning (Jacobi/Chebyshev/geometric MG) may
converge too slowly on strongly anisotropic or heterogeneous coefficients
(e.g.\ KKS mobility tensors) --- a limitation algebraic multigrid
(BoomerAMG/GAMG) is specifically built to handle but matrix-free action
alone cannot & Could force a choice between accepting slow convergence and
giving up the matrix-free zero-copy design & PR4a's mixed mode:
matrix-free \texttt{Amat} + a separately assembled approximate
\texttt{Pmat} (from \texttt{stencilAt()}) feeding
\texttt{PCHYPRE}/\texttt{PCGAMG}, preserving zero-copy on \texttt{Amat} and
the solve vectors (Section~\ref{sec:soa}) \\
R6 status & \textbf{Retired 2026-10-18} by PR1a + PR1b & Zero-copy works
on CPU and CUDA; at steady state PETSc is not slower (Jacobi-CG 3--19\%
faster; PCMG repeat solve 0.70--0.92$\times$ torch MG-PCG) & Setup cost
traced to our own configuration \\
R8 & PCMG on torch shell levels may not extend cleanly to multiple ranks:
level halos via DMDA or our own exchange, coarse gather
(\texttt{PCREDUNDANT}/\texttt{PCTELESCOPE}), coarsening of thin slabs &
Multi-rank is required; a late failure reshapes PR3 & PR2b spike \\
R9 & Non-SPD operators break Chebyshev/Jacobi + CG; torch \texttt{PCSHELL}
smoothers under GMRES are untested & SPD default does not cover all target
physics & PR2c spike \\
R10 & Fusion headroom of torch smoothers inside a PETSc callback is
unmeasured (eager V-cycle floor $\sim$2~ms in both stacks) & Decides whether
custom smoothers are worth their maintenance & PR2c, measured \\
\bottomrule
\end{longtable}

\section{Verification Methodology}
\label{sec:verification}

Every PR that touches numerics must include:

\begin{enumerate}[leftmargin=*]
\item \textbf{Method of Manufactured Solutions (MMS).} Pick a smooth
analytic $u_{\text{exact}}(x,y,t)$, substitute into the governing equation
to derive the forcing/source term, inject that source as an added
\texttt{ParsedCompute}-style term in $b$ (natural fit --- nonlinear/explicit
terms already flow into $b$ this way), run at 2--3 grid resolutions, and
confirm the solution error converges at the expected order (3-pt Laplacian:
$O(dx^2)$; 5-pt: $O(dx^4)$; first-order one-sided BC extrapolation: expect
to check whether it drops global order near boundaries --- record what's
actually observed, don't assume).
\item \textbf{Cross-validation against an independent method}, not just
regression against ourselves. Concrete anchors already in the repo:
\begin{itemize}
\item \texttt{examples/cahn\_hilliard/cahnhilliard\_fem.i} --- classic MOOSE
FEM Kernels/BCs implementation of Cahn-Hilliard. Use for PR5/PR7.
\item Marlin's own existing spectral operators
(\texttt{ReciprocalMatDiffusion}, \texttt{ReciprocalAllenCahn}) on a
periodic domain, where both the new FD path and the existing spectral path
should agree within truncation error. Use for PR3/PR5.
\item \texttt{examples/libtorch\_kks/KKS\_libtorch.i} and
\texttt{test/tests/kks/KKS\_no\_flux\_bc.i} (spectral + \texttt{psi}-mask) as
the qualitative/quantitative reference for PR6/PR7's ported KKS/SBM case.
\end{itemize}
\item \textbf{Written record attached to the PR}: convergence-order table or
plot, and the cross-validation comparison (plot or error table), not just
``tests pass.''
\end{enumerate}

\section{Development Process and Gates}
\label{sec:gates}

Every PR in the roadmap goes through the same four-stage gate before merge.
None of these stages are optional, and none can be skipped by a
plausible-sounding justification.

\begin{schematic}
\begin{verbatim}
 Implement --> Independent Reviewer --> Physics Verification --> Merge
                      Agent                     Gate
   (scope        (fresh agent/subagent,    (MMS convergence +     (status
    limited        NO shared context        cross-validation,     updated to
    to this        with implementer,        written record         "verified")
    PR's stated    runs code-review          attached to PR)
    deliverable)   skill adversarially)
\end{verbatim}
\end{schematic}

\begin{description}[leftmargin=*,style=nextline]
\item[Implement] Scope strictly to the PR's stated deliverable. No bundling
of unrelated cleanup. If something consequential is discovered
mid-implementation (e.g.\ R1 turning out false), stop and update this
document's risk/roadmap sections before continuing --- don't silently absorb
it into the PR.
\item[Independent reviewer agent] A fresh agent (not the one that wrote the
code, no inherited conversation context) runs the \texttt{/code-review}
skill on the diff, looking for correctness bugs first,
simplification/efficiency second. Findings must be fixed or explicitly
accepted with a written reason in the PR description --- ``an independent
agent reviewed this and found nothing'' is only true if one actually ran,
fresh, on the real diff.
\item[Physics verification gate] Run the Section~\ref{sec:verification}
procedure for this PR's new numerics; attach results. A PR with new
numerics and no attached convergence/cross-validation evidence does not
pass this gate, regardless of code-review outcome.
\item[Merge / go-no-go] Only after both the review and the verification
gate pass does the PR merge and its status move to \texttt{verified}. A
``go'' on a plan section is not a ``go'' on an unspecified implementation
detail discovered later --- new consequential choices get their own
sign-off.
\end{description}

\section{PR Roadmap}
\label{sec:roadmap}

Each PR is scoped to be independently reviewable and, where possible,
behavior-neutral for existing functionality until it's explicitly
activated.

\pr{0}{\texttt{LinearFDOperator} Mixin + Retrofit \texttt{FiniteDifferenceLaplacian}}
\textbf{Scope:} Add the header-only \texttt{LinearFDOperator} mixin
(\texttt{include/tensor\_computes/LinearFDOperator.h}); add it to
\texttt{FiniteDifferenceLaplacian}, implementing \texttt{stencilAt()} from
the exact same coefficients already used to build its conv kernel. No new
solver, no behavior change to any existing input file.

\textbf{Why first:} Zero-risk, establishes the interface shape everything
else builds on, and gives us a free consistency check (matrix-free kernel
vs.\ \texttt{stencilAt()} must agree) before anything depends on it.

\textbf{Verification:} Existing explicit-solve regression tests must pass
unchanged (behavior-neutral); new unit test asserts \texttt{stencilAt()}
coefficients exactly match the conv-kernel coefficients for 3-pt and 5-pt
stencil widths, 1D/2D/3D.

\textbf{Retires:} Nothing yet; sets up R4's enforcement pattern for later.

\status{proposed.}

\pr{1}{Boundary Condition TensorComputes (ghost-fill only, no PETSc yet)}
\textbf{Scope:} \texttt{TensorBoundaryCondition} base, \texttt{TensorRobinBC}
(general math), \texttt{TensorDirichletBC}/\texttt{TensorNeumannBC}
(presets); \texttt{AddTensorBCAction} parsing a new
\texttt{[TensorComputes][BCs]} sub-block; \texttt{TensorProblem::}
\texttt{addTensorComputeBC} / \texttt{\_bcs} list /
\texttt{getBoundaryConditionsFor()}; extend \texttt{exchangeGhostLayers()}
to call \texttt{fillGhostCells()} for physical (non-periodic) boundaries,
independent of partition count (fixing the current conflation of ``single
partition'' with ``no exchange needed''). \texttt{stencilContribution()} is
implemented but unused until PR4.

\textbf{Why before the solver:} Lets BCs be validated against the
\textit{existing} explicit solver (\texttt{RealSpaceForwardEuler}) first,
decoupled from PETSc risk.

\textbf{Verification:} MMS/analytic steady-state test with Dirichlet and
Neumann (zero-flux) on a small non-periodic grid, run explicitly; and the
\textbf{required R4 dual-path test} (same BC instance, check
\texttt{fillGhostCells} and \texttt{stencilContribution} imply the same
boundary condition).

\textbf{Retires:} R4, R5 (Dirichlet/Neumann usable standalone).

\status{proposed.}

\pr{1a}{Torch-Native vs.\ PETSc Matrix-Free Solve Benchmark}
\textbf{Scope:} A standalone, Python-only benchmark script (no Marlin
dependency, no MOOSE build required) comparing: (1) a pure-torch matrix-free
CG with Jacobi preconditioning, implemented directly against a
\texttt{torch.linalg.cg}-style iteration or hand-rolled CG, operating on the
same convolution-based stencil action \texttt{FiniteDifferenceLaplacian}
already uses; against (2) \texttt{petsc4py} with a \texttt{MatCreateShell}
wrapping that identical stencil action, solved via \texttt{KSPSolve} (CG +
Jacobi/Chebyshev). The \texttt{petsc4py} leg must build its \texttt{Vec}s by
wrapping the torch tensor's existing buffer zero-copy
(\texttt{VecCreateSeqWithArray}/\texttt{VecCreateMPIWithArray} on CPU;
\texttt{VecCreateSeqCUDAWithArrays}/\texttt{VecCreateMPICUDAWithArrays} or
\texttt{VecCUDAPlaceArray} on CUDA, wrapping \texttt{tensor.data\_ptr()}
directly --- see Section~\ref{sec:soa}), not a naive array copy into a
separately allocated \texttt{PETSc.Vec}; the benchmark must report both
numbers (wrapped zero-copy vs.\ naive copy) so it is clear which gap, if
any, is structural versus avoidable. Run on CPU here; the same script is
handed to the user to run on a CUDA machine (MPS is covered on this machine
for the torch leg only, since PETSc has no MPS \texttt{Vec} type to wrap at
all --- see Sections~\ref{sec:soa} and \ref{sec:decisions}). Report
wall-clock per-iteration cost, iteration count to a fixed tolerance, and the
tensor/array hand-off overhead specifically for both the wrapped and
naive-copy cases.

\textbf{Why before PR2/PR3:} This is R6 --- if torch-native matrix-free CG
is competitive with or faster than PETSc once hand-off overhead is counted,
especially on GPU, that changes how much further PETSc-specific engineering
(DMDA layout matching, options plumbing) is worth investing in before
building on it. Cheap to measure now, expensive to discover after PR2/PR3
are built.

\textbf{Verification:} The benchmark results themselves are the
deliverable; write up the comparison (CPU numbers from this machine, CUDA
numbers from the user's run) as part of this PR's record, and revisit R6's
resolution explicitly based on what's observed rather than the assumption
either way.

\textbf{Retires:} R6 (or reopens the PETSc-vs-torch-native question
explicitly if the data says so --- this PR's job is to produce real
numbers, not to confirm a predetermined answer).

\textbf{Approach actually taken:} the standalone Python/\texttt{petsc4py}
script above was abandoned in favor of a C++ unit test
(\texttt{unit/src/PetscTorchMatrixFreeBenchmarkTest.C}) in Marlin's existing
gtest harness, reusing the production \texttt{MooseTensor::}
\texttt{conjugateGradientSolve} directly instead of reimplementing CG. Two
further comparisons were added at the user's request: an anisotropic
(ill-conditioned) case testing R7's AMG claim, and a fully assembled
\texttt{Mat} scaling/memory sweep.

\textbf{Results (2026-10-17).} All correctness tests pass on both CPU
(macOS arm64, this machine) and CUDA (the user's machine, confirmed via a
second run after a device-mismatch bug found in the first run was fixed
--- see below), including the zero-copy \texttt{Vec} pointer-equality
assertion on both devices.

\textit{Matrix-free timing sweep} (CG+Jacobi, \texttt{rel\_tol}=$10^{-10}$,
unpreconditioned stopping test on both legs; 32 iterations at every size
on both legs and both devices, confirming correctness is unaffected by
which leg or device is used). CPU: macOS arm64, this machine. CUDA: H200.
CUDA times are bracketed by \texttt{cudaDeviceSynchronize()}, and
\texttt{KSPSetUp} is timed separately (0.03--0.10~ms at every size). Every
CUDA leg is solved twice back to back (``first'' and ``repeat'') for the
reason given below. The CUDA table is from run~7, which was executed under
\texttt{nsys profile} (NVTX ranges on); its absolute times are therefore
slightly inflated by the profiler (launch-bound sizes $n\le512$ by up to
$\sim$2.5~ms, large sizes by $\le$1\%, compared with the unprofiled run~5
log), so the \textit{ratios} are the result, not the millisecond values.

\begin{center}
\footnotesize
\begin{tabular}{rrrr}
\toprule
$n$ & torch-native (CPU) & PETSc (CPU, zero-copy) & PETSc / torch \\
\midrule
32 & 0.53 ms & 0.36 ms & 0.68 \\
64 & 0.94 ms & 0.71 ms & 0.76 \\
128 & 3.03 ms & 2.46 ms & 0.81 \\
256 & 11.0 ms & 8.93 ms & 0.81 \\
512 & 43.4 ms & 36.2 ms & 0.84 \\
1024 & 174 ms & 147 ms & 0.85 \\
2048 & 734 ms & 616 ms & 0.84 \\
\bottomrule
\end{tabular}
\end{center}

\begin{center}
\footnotesize
\begin{tabular}{rrrrrr}
\toprule
$n$ & torch-native & torch + D2D copy & PETSc first & PETSc repeat & PETSc repeat / torch \\
\midrule
32 & 6.68 ms & 6.74 ms & 17.5 ms & 6.07 ms & 0.91 \\
64 & 6.59 ms & 6.76 ms & 6.97 ms & 6.12 ms & 0.93 \\
128 & 6.53 ms & 6.70 ms & 7.00 ms & 6.05 ms & 0.93 \\
256 & 7.43 ms & 7.63 ms & 7.05 ms & 6.12 ms & 0.82 \\
512 & 7.05 ms & 7.24 ms & 7.94 ms & 6.66 ms & 0.95 \\
1024 & 9.08 ms & 9.18 ms & 12.4 ms & 8.79 ms & 0.97 \\
2048 & 21.5 ms & 22.0 ms & 33.3 ms & 19.9 ms & 0.93 \\
4096 & 73.7 ms & 75.8 ms & 119 ms & 65.8 ms & 0.89 \\
8192 & 278 ms & 286 ms & 463 ms & 248 ms & 0.89 \\
16384 & 1094 ms & 1128 ms & 1832 ms & 974 ms & 0.89 \\
\bottomrule
\end{tabular}
\end{center}

Torch columns are the repeat solve. The first torch solve is within 0.4\% of
it for $n\ge4096$ and within 4\% for every $64\le n\le2048$. Only $n=32$
differs substantially (21.6 against 6.7~ms), because it is the process's
first CUDA work.

\begin{center}
\includegraphics[width=\textwidth]{figures/timing_sweep.pdf}
\end{center}

\textbf{CPU.} PETSc's matrix-free shell leg is faster than torch-native at
every size: 15--19\% for $n\ge128$ (32\% at $n=32$).

\textbf{CUDA, steady state.} Below $n\approx512$ every leg sits on a
$\sim$6--7~ms floor (about 200~\textmu s per CG iteration under the
profiler), so that regime measures launch latency, not the solver. For
$n\ge4096$ the three legs are bandwidth bound ($n=16384$ is
$2.7\times10^{8}$ unknowns, 58~GB of device memory in use) and the cost is a
constant per unknown and iteration: 127~ps for torch-native, 131~ps for
torch-native with the extra copy, and 113~ps for PETSc's repeat solve.
\textbf{PETSc's repeat solve is faster than torch-native at every size,
by 3--18\% (11\% at $n\ge4096$); the small-size values are within run-to-run
noise of the launch floor and were not repeated.} It is also faster than the
copy-paying torch leg (0.80--0.96$\times$, 0.87$\times$ at $n\ge4096$), which
is the fairer comparison since \texttt{cudaShellMatMult} pays that copy too.

\textbf{First solve versus repeat solve.} PETSc's \emph{first}
\texttt{KSPSolve} after \texttt{KSPSetUp} is slower than its repeat by a
one-off cost that scales as $n^{2}$: 858~ms at $n=16384$, about 3.2~ns per
unknown at every $n\ge2048$. It is the entire apparent ``PETSc is
1.68$\times$ slower'' result of run~5; in run~7 the first-solve ratio to
torch is 1.67. Torch-native has no such cost: its repeat solve equals its
first. This was established in two steps, an \texttt{nsys} trace of run~5's
sweep and a repeat solve added in runs~6 and~7 (PETSc's repeat time is
identical in the two, 974.3~ms).

\textit{Mechanism} (read from the PETSc 3.25.4 source and matched call for
call against the trace; the \texttt{jacobi.c} lines were rechecked by hand,
but nothing here has been confirmed by an experiment yet). The cost is not in
the wrapped \texttt{b} and \texttt{x} Vecs; it is in \texttt{PCJACOBI}'s
diagonal \texttt{Vec}. \texttt{KSPSetUp} does no Jacobi work because
\texttt{jac->diag} does not exist yet, which is why it takes only 0.04~ms.
\texttt{PCApply\_Jacobi} creates the diagonal lazily on the first
application, which for the unpreconditioned stopping test happens before the
CG loop starts. \texttt{MatGetDiagonal} (the shell's callback) and
\texttt{VecReciprocal} run on the GPU. Then, because \texttt{fixdiag}
(replace a zero diagonal entry by 1) defaults on and the shell is not known
to be SPD, PETSc reads the diagonal on the host (\texttt{VecGetArray}) and
loops over all $n^{2}$ entries. That forces a pinned host allocation (the
$\sim$508~ms \texttt{cudaMallocHost}), a 39~ms device-to-host copy of the
2.1~GB vector, the host loop (the $\sim$270~ms gap, an inference: no other
code runs between the two copies), and a 39~ms copy back. This replaces my
earlier guess that PETSc round-trips a mirror of the wrapped \texttt{Vec}s;
that was wrong. No CG iteration runs before the stall.

\textit{Recurrence.} From the source: not for the same \texttt{KSP} with
unchanged operators, including after wrapping new \texttt{Vec}s over the
same buffers. It does recur in full for a newly created \texttt{KSP}/\texttt{PC}. If
the operator's state changes on an existing \texttt{PC}, the diagonal is
rebuilt and the host loop and both copies return, but not the pinned
allocation. Whether the Marlin Picard loop changes the \texttt{Mat} state
every iteration is not yet known and decides how often this is paid.

The cost belongs to \texttt{PCJACOBI}'s setup rather than to zero-copy
wrapping, so a different preconditioner would not pay it. It is avoidable
by configuration, but each option is a numerical choice:
\texttt{PCJacobiSetFixDiagonal(pc, PETSC\_FALSE)} would stop zero diagonal
entries being replaced by 1 (identical results here, where the diagonal is a
nonzero constant) and \texttt{MatSetOption(mat, MAT\_SPD, PETSC\_TRUE)}
asserts the operator is SPD. Neither has been tried. At 3.2~ns per unknown
the one-off cost is about 80\% of an entire 32-iteration torch-native solve
(4.1~ns per unknown), so creating the \texttt{KSP}/\texttt{PC} once and
reusing it across solves matters until this is settled. PR3 must verify
this on its own code path.

\textbf{Where PETSc's steady-state advantage comes from.} The matvec is the
same work in every leg: the \texttt{im2col} and \texttt{gemv} kernels total
501~ms of GPU time at $n=16384$ in each leg (46\% of torch-native's busy
time). Everything else is vector work. Torch-native spends 582~ms there in
many ATen elementwise and reduction kernels; PETSc spends about 499~ms in
cuBLAS \texttt{axpy}, \texttt{dot} and \texttt{nrm2} (first-solve trace,
which also contains the 78~ms of host copies counted separately, so this
is only a rough comparison). That fits the observed 14~ps, or $\sim$120~ms
per solve, difference, but it is a reading of kernel totals, not a
controlled experiment. Neither leg avoids synchronization: both perform
three blocking scalar reductions per CG iteration (about 3 device-to-host
copies per iteration in each leg's trace), so synchronization does not
explain the difference, and eager libtorch fuses nothing. The extra
device-to-device copy in \texttt{cudaShellMatMult} costs 3.9~ps per unknown
and iteration (3.1\%), as the copy leg shows directly.

R6 therefore resolves in PETSc's favour on both devices at steady state:
zero-copy \texttt{Vec} wrapping (Section~\ref{sec:soa}) works correctly on
CPU and CUDA, and the matrix-free PETSc leg is faster than torch-native on
both, by 15--19\% on CPU and 3--18\% on this GPU (11\% at large $n$). The condition is the
reuse requirement above; the first-solve cost is the open item.

\textit{Ill-conditioned (anisotropic, R7/PR4a) comparison} ($n=64$,
$D_x=1$, $D_y=10^{-4}$, $D_x/D_y=10^4$, same manufactured solution and
correctness check in both legs). \texttt{KSPSetUp} is called before the
timer so that PETSc's lazy \texttt{PCSetUp} (the GAMG hierarchy build) is
not charged to the first solve; ``repeat solve'' is a second solve from a
zero guess on the same \texttt{KSP} with the preconditioner reused, with the
iteration count and solution error asserted equal to the first solve.

\begin{center}
\footnotesize
\begin{tabular}{lrrrrr}
\toprule
Case & Its. & \texttt{Pmat} assembly & PC setup & First solve & Repeat solve \\
\midrule
Jacobi (CPU) & 202 & --- & 0.01 ms & 4.37 ms & 4.34 ms \\
AMG (CPU) & 134 & 0.33 ms & 3.77 ms & 18.4 ms & 18.2 ms \\
Jacobi (CUDA) & 202 & --- & 0.06 ms & 27.7 ms & 26.7 ms \\
AMG (CUDA) & 134 & 0.67 ms & 6.94 ms & 80.1 ms & 79.2 ms \\
\bottomrule
\end{tabular}
\end{center}

\begin{center}
\includegraphics[width=\textwidth]{figures/anisotropic_amg_vs_jacobi.pdf}
\end{center}

AMG needs 1.5$\times$ fewer iterations (134 against 202), identically on
both devices. That is concrete evidence for R7 that a scalar matrix-free
preconditioner cannot reach on its own. It does not translate into
wall-clock time at this size. With the preconditioner already built, an AMG
solve is 4.2$\times$ slower than a Jacobi solve on CPU (18.2 against
4.3~ms) and 3.0$\times$ slower on CUDA (79.2 against 26.7~ms). The one-time
costs (\texttt{Pmat} assembly plus PC setup) are only 4.1~ms on CPU and
7.6~ms on CUDA, so amortizing them over many solves cannot rescue AMG
here; the cost is in applying the hierarchy each iteration. Two further
points. First, $n=64$ (4096 unknowns) is deep in the CUDA latency-floor
regime: 26.7~ms over 202 iterations is again $\sim$130~\textmu s per
iteration, so the CUDA anisotropic numbers say little about AMG on a GPU.
Second, \texttt{Pmat} is a host \texttt{MATSEQAIJ}, so GAMG's hierarchy and
V-cycle run on the host and each \texttt{PCApply} presumably moves the
residual and correction across the bus. That would explain AMG being 4.3$\times$
slower on CUDA than on CPU (79.2 against 18.2~ms), but it has not been
measured. Whether AMG pays off at larger $n$ (where Jacobi-CG iteration
counts are expected to grow and AMG's to stay roughly flat) is untested.
Whether \texttt{Pmat} is built once or per step depends on the reassembly
cadence of Section~\ref{sec:decisions} item 3, whose v1 default is every
Picard iteration; at these costs that choice does not change the
conclusion.

\textbf{History of this comparison.} The first two CUDA runs reported
identical Jacobi/AMG iteration counts to the CPU run because the test
(\texttt{AnisotropicAMGConvergesFasterThanJacobi}) was CPU-only; a genuine
CUDA variant was added (commit \texttt{cb9368f}). Wall-clock timing was
added next (commit \texttt{24527aa}) because iteration counts alone could
not show which leg is faster. The initial wall-clock numbers charged
PETSc's lazy \texttt{PCSetUp} to the first solve; the setup/solve/repeat
split above (run~5) removes that. The CUDA matrix-free sweep was
originally reported as ``PETSc is 1.68$\times$ slower, $\sim$82~ps per unknown
and iteration unexplained''. That compared PETSc's first solve with
torch-native's first solve, and the gap was PETSc's one-off cost, not
per-iteration work. Adding repeat solves to every leg (runs~6 and~7)
withdrew that finding; the $\sim$82~ps figure and the plateau claim no
longer stand.

\textit{Assembled-\texttt{Mat} fallback scaling} (same isotropic problem as
the timing sweep, a fully assembled \texttt{Mat} for both \texttt{Amat} and
the Jacobi \texttt{Pmat}, CPU only --- not yet exercised on CUDA):

\begin{center}
\footnotesize
\begin{tabular}{rrrr}
\toprule
$n$ & assembly & solve & \texttt{Mat} storage (\texttt{nz\_used}) \\
\midrule
32 & 0.08 ms & 0.09 ms & 0 MB \\
256 & 4.32 ms & 5.44 ms & 5 MB \\
1024 & 67.8 ms & 94.6 ms & 80 MB \\
2048 & 276 ms & 394 ms & 320 MB \\
\bottomrule
\end{tabular}
\end{center}

Iteration counts match the matrix-free legs exactly at every size, so
assembly changes nothing about correctness or convergence, only cost. At
$n=2048$ the assembled matrix occupies 320~MB on its own. Once assembled,
its solve (394~ms) is faster than either matrix-free CPU leg (616 and
734~ms); with assembly included (670~ms) it falls between them. For this
constant-coefficient 5-point stencil on CPU, the evidence for the
matrix-free default (Section~\ref{sec:decisions}, item 5) is therefore
storage and device residency, not speed. \textbf{Peak RSS is deliberately
not tabulated:} \texttt{ru\_maxrss} is a process-wide high-water mark and
this sweep runs after the matrix-free sweep in the same process (which had
already reached 1121~MB), so the log's RSS column does not measure the
assembled path's own footprint. An earlier revision of this section
compared those numbers; that comparison was not valid and is withdrawn.

\textbf{Fix found and applied mid-verification:} the first CUDA run
(before the fixes above) failed both new CUDA-only tests with a device
mismatch (\texttt{conv2d}'s weight tensor was unconditionally CPU-resident
in \texttt{PeriodicStencilOperator}, so applying it to a CUDA input threw).
Fixed by giving the stencil operator classes an optional \texttt{device}
parameter (commit \texttt{19341d7}); a second CUDA run with the fix passed
all 8 tests. This is recorded here, rather than only in commit history,
because it is exactly the kind of PETSc/torch hand-off correctness issue
Section~\ref{sec:soa}'s zero-copy requirement is meant to get right, and it
only surfaced once real CUDA hardware exercised the code.

The full raw logs from both machines are reproduced below for direct
inspection, rather than only the extracted tables above.

\lstinputlisting[caption={Full CPU run log (\texttt{petsc\_torch\_benchmark\_cpu.log}), run~5 (setup/solve/repeat timing), 6/6 tests passing.}]{logs/petsc_torch_benchmark_cpu.log}

\lstinputlisting[caption={Full CUDA run log (\texttt{petsc\_torch\_benchmark\_cuda.log}), run~5 on an H200, all 9/9 tests passing. Its \texttt{DISABLED\_CudaTimingSweep} lines time first solves only and are superseded by the repeat-solve log below; the other tests' lines (including the anisotropic CUDA test) are current.}]{logs/petsc_torch_benchmark_cuda.log}

\lstinputlisting[caption={CUDA timing sweep with first and repeat solves for every leg (\texttt{petsc\_torch\_benchmark\_cuda\_sweep.log}), run~7 on an H200 under \texttt{nsys profile}, 1/1 tests passing.}]{logs/petsc_torch_benchmark_cuda_sweep.log}

\status{implemented and verified on both CPU and CUDA, all 9 tests,
2026-10-17 (commits \texttt{12016ae}, \texttt{41c2e31}, \texttt{19341d7},
\texttt{cb9368f}, \texttt{24527aa}, \texttt{c77191a}, \texttt{5c38c75},
\texttt{d0f8d60}, \texttt{3ac2c7c} on branch \texttt{petsc\_torch\_benchmark}).
R6 and R7 are both evidenced on CPU and CUDA. At steady state PETSc's
matrix-free leg is faster than torch-native on both devices (CUDA: 11\% at
large $n$). Its first \texttt{KSPSolve} carries a one-off, $n^{2}$-scaling
cost (3.2~ns per unknown) traced from the PETSc source to
\texttt{PCJACOBI}'s lazily built, host-checked diagonal (not yet confirmed by
an experiment); PR3 must create the \texttt{KSP}/\texttt{PC} once and reuse it.
Remaining before this PR can be marked \texttt{verified} in the
status-tracking table: the independent-reviewer-agent pass
(Section~\ref{sec:gates}) has not yet been run on this diff.}

\pr{1b}{Multigrid Spikes and the Hybrid-Design Decision}
\label{sec:pr1b}
These are spikes, not maintained solvers, on branch \texttt{petsc\_torch\_benchmark}.
Unless stated otherwise every run is 2D, periodic, with the SPD operator
$-\Delta u + \sigma u$ ($\sigma = 1$). The CUDA runs were done by the user on an
H200 with PETSc 3.25.4 and libtorch 2.10. The table lists the tests in the
order they were run.

\begin{longtable}{>{\bfseries}p{0.03\textwidth} p{0.25\textwidth} p{0.4\textwidth} p{0.23\textwidth}}
\toprule
\# & Test (commit) & Result & Led to \\
\midrule
\endhead
1 & PR1a Jacobi-CG benchmark, CPU + CUDA (\texttt{12016ae}\ldots\texttt{3ac2c7c}) & Zero-copy holds on both devices. At steady state PETSc is 15--19\% faster on CPU and 3--18\% faster on CUDA. The first CUDA solve is slow & Why? \\
2 & First-solve source trace (\texttt{27a8807}) & \texttt{PCJACOBI} lazily builds its diagonal and checks it on the host; the \texttt{Vec} wrapping is not the cause & Jacobi-CG is $O(n)$ iterations: is torch GMG feasible? \\
3 & Torch GMG, emulated slabs (\texttt{c3364ec}) & RBGS V-cycle MG-PCG takes 11 iterations at $n=64..512$, identical bit for bit on 1/2/4 slabs & Real MPI? \\
4 & Torch GMG, real MPI (\texttt{34e10ff}) & 10 iterations at np $=1/2/4/8$ with the same error. Dropping one halo update breaks it & DMDA instead? \\
5 & DMDA decomposition spike (\texttt{18d5886}, unmerged) & A DMDA forced onto Marlin's partition matches ownership and ghosts and aliases the data zero-copy on 1--8 ranks & R1 feasible \\
6 & Torch V-cycle CUDA sweep (\texttt{e566e3c}) & $\sim$2~ms eager launch floor. MG-PCG takes 11 iterations at every $n$ and beats Jacobi-CG by 4.5$\times$ to $>$360$\times$ & Does PETSc GMG match? \\
7 & PETSc KSPCG+PCMG on torch shells, CPU (\texttt{610802e}) & 9 iterations, grid independent & Time it on the GPU \\
8 & PCMG CUDA sweep (\texttt{6698261}) & 8--9 iterations. Repeat solve is 0.70--0.92$\times$ torch's, with the same V-cycle floor. \texttt{KSPSetUp} costs 1.5--1.8$\times$ a solve & Paid every Picard step? \\
9 & Picard update-policy sweep (\texttt{bd00ff4}, \texttt{665ac06}, \texttt{be91418}) & A state bump triggers a lazy rebuild. Stale setups can go indefinite, so retry is mandatory. Threshold reuse saves 39\% at $\beta=3$. At $\beta=9$ every policy fails the DTOL test ($\|b\|$ reference with a warm start) & Stopped tuning the test; audited our PETSc usage \\
10 & \texttt{-log\_view} audit & Every host$\leftrightarrow$device copy is in \texttt{PCSetUp}: the Jacobi zero-diagonal check and Chebyshev's GMRES estimate. \texttt{PCApply} and \texttt{MatMult} make no copies & Tell PETSc the operator is SPD \\
11 & \texttt{MAT\_SPD} variant (\texttt{2ded6f0}) & Setup 724$\to$408~ms at $n=8192$ and 69$\to$9~ms at $n=256$; nothing else changes. The Chebyshev estimate remains & Gershgorin bounds deferred; hybrid design adopted \\
\bottomrule
\end{longtable}

Logs: \texttt{doc/roadmap/logs/torch\_multigrid\_cuda\_vcycle\_sweep.log},
\texttt{petsc\_multigrid\_cuda\_sweep.log} and
\texttt{petsc\_multigrid\_cuda\_sweep\_log\_view.log} (gitignored).

\status{implemented. CPU tests pass and the CUDA sweeps have been run. No
independent review yet. These are spikes only and are not on the production
path.}

\pr{2}{DMDA/Decomposition Compatibility Spike}
\textbf{Scope:} Isolated, throwaway-if-needed investigation + the minimal
\texttt{DomainAction} addition (\texttt{getGlobalIndexOffset()} or
equivalent) needed to build DMDA's explicit \texttt{lx/ly/lz} partition
arrays. Deliverable is a small standalone test program/test case that
builds a DMDA matching Marlin's existing decomposition on 1, 2, and
$>$2 MPI ranks and verifies, by direct index comparison, that DMDA-owned
ranges match Marlin's torch-tensor local ownership exactly --- no PDE solve
yet.

\textbf{Why isolated:} This is R1, the highest-risk unknown. If DMDA can't
be forced to match, the whole PETSc-wiring design in PR3+ needs to change
(fallback: hand-rolled DOF numbering) --- better to find out now than
inside a bigger PR.

\textbf{Verification:} The index-matching test itself \textit{is} the
verification; report pass/fail across 1/2/4+ rank configurations and both
periodic and non-periodic axes.

\textbf{Retires:} R1.

\status{spike implemented (\texttt{18d5886}, agent worktree, unmerged). It
passes on 1/2/3/4/6/8 ranks using only public \texttt{DomainAction} getters.
Non-periodic axes and review are outstanding.}

\pr{2b}{Multi-Rank PCMG on Torch Shell Levels (Spike)}
\textbf{Scope:} Run PR1b's PCMG problem on real MPI ranks through
\texttt{DomainAction}. Settle the level halo exchange (per-level DMDA
vs.\ Marlin's own exchange), the coarse gather
(\texttt{PCREDUNDANT}/\texttt{PCTELESCOPE}), and how to coarsen thin slabs.
Add a one-rank 3D check.

\textbf{Verification:} Iteration count and solution must be independent of
the rank count (np $=1/2/4/8$). Dropping one halo update must break
convergence.

\textbf{Retires:} R8. \status{proposed.}

\pr{2c}{Non-SPD Operator with a Torch \texttt{PCSHELL} Smoother + Fusion Measurement (Spike)}
\textbf{Scope:} Advection--diffusion solved with KSPGMRES/FGMRES + PCMG.
Compare a torch symmetric red-black Gauss-Seidel \texttt{PCSHELL} smoother
against PETSc's own smoother. Time the same smoother eager, fused through
\texttt{TracedComputeSequence}, and CUDA-graph captured.

\textbf{Verification:} Iterations and solution against a reference solve;
timings reported, not assumed.

\textbf{Retires:} R9, R10. \status{proposed.}

\pr{3}{\texttt{PetscFDSolverBase}, Single Field, Matrix-Free, No Picard, No Nonlinear Remainder}
\textbf{Hybrid-design update (2026-10-18):}
\begin{itemize}[leftmargin=*]
\item \textbf{Default solver for an SPD field.} KSPCG + PCMG with torch shell
level operators (rediscretised on each level), torch shell transfers,
Chebyshev/Jacobi smoothers and an assembled coarse \texttt{Mat}. All of this
can be overridden through \texttt{-<field>\_} options.
\item \textbf{SPD marking.} Levels are marked \texttt{MAT\_SPD} when the
operator declares itself SPD.
\item \textbf{Chebyshev bounds.} Leave a hook for analytic Chebyshev bounds.
\item \textbf{DTOL with warm starts.} PETSc's default divergence test uses
$\|b\|$ as its reference when the initial guess is nonzero. Whether to
change it is a separate design decision, not a default change.
\item \textbf{Extra verification.} \texttt{-log\_view} must show no
host$\leftrightarrow$device copies in \texttt{KSPSolve}, and the iteration
count must be independent of grid size.
\end{itemize}

\textbf{Scope:} \texttt{PetscFDSolverBase : public TensorSolver} --- the
single solver object, exercised here with exactly one tracked field and one
\texttt{LinearFDOperator} contribution (constant-D diffusion or Laplacian),
no boundary terms yet (periodic domain only), reusing
\texttt{Moose::PetscSupport} for \texttt{-<field>\_} prefixed options. The
field's \texttt{LinearFDOperator} action is bound into a
\texttt{MatCreateShell} whose \texttt{MatMult} callback invokes the existing
matrix-free conv-kernel code directly on the torch tensor --- no sparse
\texttt{Mat} is assembled. The field's solution/RHS \texttt{Vec}s are built
by wrapping its existing torch tensor buffer zero-copy
(Section~\ref{sec:soa}: \texttt{VecCreateMPIWithArray} on CPU,
\texttt{VecCreateMPICUDAWithArrays}/\texttt{VecCUDAPlaceArray} on CUDA), not
by allocating separate PETSc-owned storage. \texttt{substep()}'s
one-field-at-a-time loop degenerates to binding the shell operator + one
\texttt{KSPSolve} when there's only one tracked field; the loop structure
itself (so later PRs just add fields, not new solver plumbing) is part of
this PR's deliverable.

\textbf{Verification:} MMS convergence-order study on a periodic domain;
cross-validation against the existing spectral
\texttt{ReciprocalLaplacianFactor}/\texttt{FFTSemiImplicit} path on the same
problem; confirm \texttt{petsc\_options\_iname}/\texttt{\_value} set in the
input file actually reach the KSP (e.g.\ via \texttt{-<field>\_ksp\_view}
showing the configured PC type and confirming no AMG/assembled-matrix
option is silently required). Confirm zero-copy directly: compare
\texttt{data\_ptr()} of the torch tensor against the address PETSc's
\texttt{Vec} resolves to, and/or a profiler trace showing no
\texttt{cudaMemcpy} during \texttt{KSPSolve} on the CUDA build --- not
inferred from timing alone. Profile matrix-free action + solve cost (R3,
first measurement) and compare directly against PR1a's standalone benchmark
numbers on the same problem size, using PR1a's wrapped-zero-copy numbers as
the relevant comparison.

\textbf{Retires:} Nothing new, but is the first real exercise of R1's fix
and the PETSc-options-reuse mechanism end to end, in the matrix-free mode
R6/PR1a is expected to favor.

\status{proposed.}

\pr{4}{Boundary Terms in Matrix-Free Operator Application (connect PR1 + PR3)}
\textbf{Scope:} \texttt{PetscFDSolverBase}'s \texttt{MatMult} callback and
\texttt{assembleRHS} consult \texttt{TensorProblem::getBoundaryConditionsFor()}
near physical edges, calling a matrix-free
\texttt{applyGhostContribution()}-style method on
\texttt{TensorBoundaryCondition} (the matrix-free analog of what was
originally named \texttt{stencilContribution()}); run the single-field
implicit solver on a non-periodic domain with Dirichlet/Neumann/Robin
boundaries.

\textbf{Verification:} MMS with a non-trivial boundary term; confirm
observed convergence order near boundaries (record actual order, don't
assume); the R4 dual-path test from PR1 is re-run in this matrix-free
implicit context too (not just the explicit context).

\textbf{Retires:} Closes the loop on R4 across the explicit and matrix-free
implicit paths.

\status{proposed.}

\pr{4a (build if PR4's matrix-free preconditioning proves insufficient on anisotropic/heterogeneous cases)}{AMG-Preconditioned Mixed Mode}
\textbf{Scope:} A \texttt{Pmat} assembly path
(\texttt{assemblePreconditionerMatrix()}) building a real sparse PETSc
\texttt{Mat} via \texttt{MatSetValuesStencil} from the same
\texttt{LinearFDOperator} coefficients (exposed via \texttt{stencilAt()}),
passed to \texttt{KSPSetOperators(ksp, Amat, Pmat)} as the second argument,
alongside the existing matrix-free \texttt{Amat} (\texttt{MatCreateShell}),
which remains in place for the Krylov iteration itself.
\texttt{-<field>\_pc\_type hypre}/\texttt{gamg} then builds BoomerAMG/GAMG
from \texttt{Pmat} while every \texttt{MatMult} still runs against the
zero-copy, matrix-free \texttt{Amat} (Section~\ref{sec:soa}). \texttt{Pmat}
is reassembled at the preconditioner-update cadence of
Section~\ref{sec:decisions} item 3 (PCMG on torch shells is now PR3's
default; PR4a remains the route where geometric coarsening fails); it may be an
approximate/frozen-coefficient matrix rather than an exact Jacobian, since
AMG preconditioners only need the right coarsening structure. Selected
per-field via a solver option. A separate, lower-priority fallback ---
assembling \texttt{Amat} itself as a real sparse \texttt{Mat} instead of a
shell, forfeiting zero-copy entirely --- is kept as a documented last
resort, scoped only if the mixed mode itself proves insufficient.

\textbf{Rationale:} This is R7. Strongly anisotropic or heterogeneous
coefficient fields (plausible in KKS/grand-potential mobility tensors, not
just hypothetical) are a case matrix-free Jacobi/Chebyshev/geometric-MG
preconditioning is known to converge slowly on, and algebraic multigrid is
specifically built to handle via coarsening on actual matrix entries. The
mixed-\texttt{Pmat} mode gets AMG's convergence benefit without
reintroducing the per-solve \texttt{Vec}/\texttt{Mat} hand-off cost
Section~\ref{sec:soa}'s zero-copy requirement exists to avoid --- only
\texttt{Pmat}'s own assembly is paid, at the Section~\ref{sec:decisions}
item 3 cadence, not every solve.

\textbf{Verification:} Side-by-side iteration-count/wall-clock comparison
against PR4's pure matrix-free path on an anisotropic test case, to confirm
the convergence benefit materializes and quantify \texttt{Pmat} assembly
cost against the solve cost it buys.

\status{proposed (contingent --- may never be built if PR4's pure
matrix-free preconditioning suffices on the target physics).}

\pr{5}{\texttt{FiniteDifferenceDiffusion} (variable coefficient) + lagged nonlinear remainder in $b$}
\textbf{Scope:} New real-space variable-coefficient $\nabla\cdot(D\nabla u)$
operator implementing \texttt{LinearFDOperator} (harmonic-mean half-point
coefficients); wire a generic lagged-nonlinear-remainder path into
\texttt{assembleRHS} so a non-\texttt{LinearFDOperator} TensorCompute's
\texttt{computeBuffer()} output contributes to $b$ at the current/lagged
state.

\textbf{Verification:} MMS with spatially varying $D$; cross-validation
against \texttt{examples/cahn\_hilliard/cahnhilliard\_fem.i} (independent
FEM implementation) and/or the existing spectral
\texttt{ReciprocalMatDiffusion} on a periodic case. This is the first PR
that actually exercises a full
Cahn-Hilliard-shaped single-field implicit-diffusion / explicit-nonlinear
split.

\textbf{Retires:} Nothing new; first real profiling data point for R3 under
a nontrivial coefficient field.

\status{proposed.}

\pr{5a}{\texttt{TimeDerivativeOperator} + Steady-State Field Mixing (Poisson-alongside-phase-field)}
\textbf{Scope:} New \texttt{TimeDerivativeOperator}
(\texttt{LinearFDOperator}-implementing) making the time-derivative term an
explicit, optional per-field contribution rather than an assumption baked
into \texttt{PetscFDSolverBase} (see Section~\ref{sec:steady-mixing});
demonstrate a steady field (e.g.\ a Poisson equation for a potential) with
no \texttt{TimeDerivativeOperator} registered, tracked in the same
\texttt{PetscFDSolverBase} one-by-one field loop alongside a transient
field, with a lagged source term coupling the two.

\textbf{Verification:} MMS for the steady Poisson solve in isolation
(confirm it reduces to a standard steady diffusion MMS case with no time
stepping); a coupled two-field case (one steady, one transient) checking
the steady field is being re-solved to a fresh steady state each Picard
iteration/timestep rather than going stale.

\textbf{Retires:} Closes out the user's steady+transient mixing
requirement.

\status{proposed.}

\pr{6}{\texttt{PicardTensorSolver} + Segregated Two-Field Coupling (KKS-shaped)}
\textbf{Scope:} \texttt{PicardTensorSolver : public TensorSolver} wrapping a
\texttt{PetscFDSolverBase} with $\geq 2$ tracked fields (e.g.\ $c, \eta$),
reusing the \texttt{IterativeTensorSolverInterface} knob shape (min/max its,
rel/abs tol or Postprocessor-driven convergence, relaxation factor).

\textbf{Verification:} Measure actual Picard convergence behavior (this is
R2 --- do not assume it converges; report iteration counts and residual
decay for the real coupled case, tune relaxation if needed); cross-validate
against \texttt{examples/libtorch\_kks/KKS\_libtorch.i}'s qualitative
behavior; re-profile per-iteration cost under the item-3 reuse policy
across two coupled fields (R3, second measurement).

\textbf{Retires:} R2, R3 (final measurement before declaring the
performant-goal met or needing follow-up work).

\status{proposed.}

\pr{7}{Port SBM/KKS No-Flux Example to the New Pipeline}
\textbf{Scope:} Reimplement \texttt{test/tests/kks/KKS\_no\_flux\_bc.i}'s
physics using the new \texttt{TensorNeumannBC} (zero-flux) in place of the
\texttt{psi}-masking hack, on the new FD+PETSc solve path.

\textbf{Verification:} Compare interface evolution/qualitative behavior
against the original spectral \texttt{psi}-masked version; conservation
checks (e.g.\ total concentration) over the run; this is the end-to-end
demonstration that the ``basic'' target functionality (KKS + SBM) is met.

\textbf{Retires:} Closes out the stated ``basic'' functionality goal.

\status{proposed.}

\pr{8 (deferred)}{Monolithic Block-Matrix Groundwork}
\textbf{Scope:} Not scheduled; documented here so the \texttt{target\_field}
tagging convention from PR0 doesn't rot unused. Revisit only if a future
model needs exact (non-Picard) multi-field coupling.

\status{deferred --- explicitly out of scope until requested.}

\section{Status Tracking}

\begin{longtable}{>{\bfseries}p{0.08\textwidth} p{0.42\textwidth} p{0.18\textwidth} p{0.2\textwidth}}
\toprule
PR & Deliverable & Status & Verification attached? \\
\midrule
\endhead
PR0 & LinearFDOperator mixin + Laplacian retrofit & proposed & --- \\
PR1 & BC TensorComputes, ghost-fill & proposed & --- \\
PR1a & Torch-native vs.\ PETSc matrix-free benchmark & implemented, CPU+CUDA verified & CPU+CUDA logs and tables in Section~\ref{sec:roadmap} (PR1a) \\
PR1b & Multigrid spikes $\to$ hybrid design & implemented; review not run & chain table and logs (Section~\ref{sec:pr1b}) \\
PR2 & DMDA/decomposition compatibility spike & spike implemented, unmerged & passes 1--8 ranks \\
PR2b & Multi-rank PCMG on torch shells & proposed & --- \\
PR2c & Non-SPD + torch PCSHELL smoother + fusion & proposed & --- \\
PR3 & PetscFDSolverBase, single field, matrix-free & proposed & --- \\
PR4 & BC terms in matrix-free operator application & proposed & --- \\
PR4a & AMG-preconditioned mixed mode (matrix-free Amat + assembled Pmat), contingent & proposed & --- \\
PR5 & FiniteDifferenceDiffusion + lagged remainder & proposed & --- \\
PR5a & TimeDerivativeOperator + steady-field mixing & proposed & --- \\
PR6 & PicardTensorSolver, two-field coupling & proposed & --- \\
PR7 & SBM/KKS no-flux port & proposed & --- \\
PR8 & Monolithic groundwork & deferred & --- \\
\bottomrule
\end{longtable}

\section{Open Questions Deferred (not blocking PR0)}

\begin{enumerate}[leftmargin=*]
\item \sout{Input syntax for multiple PETSc solvers} --- \textbf{resolved:}
a single \texttt{PetscFDSolverBase} instance per \texttt{[TensorSolver]}
block lists its tracked fields and iterates over them one-by-one; each
field's \texttt{petsc\_options*} triplet is a named sub-block keyed by field
name within that one \texttt{[TensorSolver]} block. No separate per-field
solver objects.
\item DMDA stencil width when different linear operators need different
ghost widths --- v1 sizes one DM to the max; revisit only if that proves
wasteful in PR3/PR5 profiling.
\item \sout{Whether \texttt{coefficientsChanged()} caching is worth
enabling} --- superseded by decision \#3's reuse policy. Open: the threshold
value and whether it is user-facing.
\item Whether DMDA is still the right choice (PETSc wiring design) once
matrix-free is the default --- DMDA's main unique value was correct
assembled-matrix DOF numbering; under matrix-free action, Marlin's existing
MPI decomposition may be sufficient on its own without a DMDA layer at all
(PETSc's \texttt{Vec} layout still needs to match for \texttt{KSPSolve}'s
own vector operations, but that's a smaller ask than matching \texttt{Mat}
row ownership). Revisit when scoping PR2 in light of PR1a's results.
\item Whether PETSc is used at all for the matrix-free path, or whether
PR1a's benchmark shows a pure torch-native CG/multigrid solver (no PETSc
dependency) is preferable --- this is R6, and is now a per-device question
rather than one global answer: on CPU and CUDA, zero-copy \texttt{Vec}
wrapping (Section~\ref{sec:soa}) is structurally possible, so the question
is whether it is faster in practice once wrapped (PR1a decides this with
real numbers); on MPS, no PETSc \texttt{Vec} type exists at all, so the
torch-native path is the only option there regardless of what PR1a
measures on CPU/CUDA. \textbf{Resolved 2026-10-18:} PETSc (KSPCG+PCMG on
torch shells) is the default on CPU and CUDA (Section~\ref{sec:hybrid}).
\item Which solves PR3--PR6 need first (SPD-only vs.\ non-SPD), and whether
MPS is a first-class target or torch-CG-only.
\item Whether non-periodic BCs need special treatment on coarse PCMG levels
(rediscretised BC rows vs.\ Galerkin) --- check in PR2b/PR4.
\end{enumerate}

\end{document}
